\newcommand{\verfasser}{Phil Steinhorst}					% Verfasser des Dokuments
\newcommand{\fach}{Elementare Zahlentheorie}				% Titel der Vorlesung								
\newcommand{\shortFach}{EZT}								% Titel kurz
\newcommand{\prof}{Prof.\,Dr.\,Falko Lorenz}					% Dozent
\newcommand{\untertitel}{gelesen von \prof}					% Untertitel
\newcommand{\semester}{Wintersemester 2014/2015}					% Semester
\newcommand{\homepage}{http://wwwmath.uni-muenster.de/u/karin.halupczok/elZTWiSe14/}	% Vorlesungshomepage
% % % % % % % % % % % % % % % % % % % % % % % % % % % % % % % % %
%		PhistScript.tex											%
%		XeLaTeX-Settings for Scripts							%
%																%
%		Author: Phil Steinhorst									%
% % % % % % % % % % % % % % % % % % % % % % % % % % % % % % % % %

\setlength\parindent{0pt}             % ohne Einrueckung

% % % Konfiguration von scrheadings
\setheadsepline{1pt}[\color{black}]
\pagestyle{scrheadings}
\clearscrheadfoot

% % % Kopf-/Fußzeilenlayout
\providecommand{\shortFach}{\fach}
\ohead{\small \leftmark}
\ofoot[{ \small \thepage}]{ \small \thepage} 
\automark{section}

% Metadaten
\title{\fach}
\author{\verfasser}
\subtitle{\untertitel}
\date{\semester}

% Inhaltsverzeichnis
\setcounter{tocdepth}{1}
\usepackage[tocindentauto]{tocstyle}
\usetocstyle{KOMAlike}	

% Nummerierung
\numberwithin{equation}{section}		% Section in Equation-Nummerierung einbeziehen

% ntheorem package
\newcounter{tcount}						% Zähler für Theoreme
\numberwithin{tcount}{section}			% Section in Theorem-Nummerierung einbeziehen
\usepackage[hyperref]{ntheorem}			% Lade ntheorem mit hyperref-Workaround
\theoremstyle{break}					% Zeilenumbruch nach Theoremkopf
\theorembodyfont{\normalfont}			% nicht kursiv
\theorempreskipamount0.5cm				% Abstand vor Theorem
\theorempostskipamount0.5cm				% Abstand nach Theorem
\newtheorem{thm}[tcount]{Theorem}
\newtheorem{bsp}[tcount]{Beispiel}
\newtheorem{defn}[tcount]{Definition}
\newtheorem{bem}[tcount]{Bemerkung}
\newtheorem{lemma}[tcount]{Lemma}

% Workaround für Tabulatoren im Theorem-Verzeichnis
\makeatletter
\def\thm@@thmline@name#1#2#3#4{%
        \@dottedtocline{-2}{0em}{2.3em}%
                   {\makebox[\widesttheorem][l]{#1 \protect\numberline{#2}}#3}%
                   {#4}}
\@ifpackageloaded{hyperref}{
\def\thm@@thmline@name#1#2#3#4#5{%
    \ifx\\#5\\%
        \@dottedtocline{-2}{0em}{2.3em}%
            {\makebox[\widesttheorem][l]{#1 \protect\numberline{#2}}#3}%
            {#4}
    \else
        \ifHy@linktocpage\relax\relax
            \@dottedtocline{-2}{0em}{2.3em}%
                {\makebox[\widesttheorem][l]{#1 \protect\numberline{#2}}#3}%
                {\hyper@linkstart{link}{#5}{#4}\hyper@linkend}%
        \else
            \@dottedtocline{-2}{0em}{2.3em}%
                {\hyper@linkstart{link}{#5}%
                  {\makebox[\widesttheorem][l]{#1 \protect\numberline{#2}}#3}\hyper@linkend}%
                    {#4}%
        \fi
    \fi}
}
\makeatother
\newlength\widesttheorem
\AtBeginDocument{
  \settowidth{\widesttheorem}{Proposition 10.10\quad}	% Referenzstring für die Breite der ersten Spalte
}

\theoremlisttype{optname}	% nur Theoreme auflisten, die einen Namen haben. 							% XeLaTeX-Konfiguration für Skripte

% % % Dokumentspezifische Einstellungen % % %
\setcounter{tocdepth}{2}				% Tiefe im Inhaltsverzeichnis (1: nur Sections, 2: auch subsections...)

\numberwithin{equation}{section}			% Section in Equation-Nummerierung einbeziehen
\newcounter{countsatz}						% Zähler für Sätze
\numberwithin{countsatz}{section}			% Section in Satz-Nummerierung einbeziehen
\newcounter{countfalko}						% Zähler für Fs
\numberwithin{countfalko}{section}			% Section in F-Nummerierung einbeziehen
\newcounter{countbsp}						% Zähler für Beispiele
\numberwithin{countbsp}{section}			% Section in Beispiel-Nummerierung einbeziehen
\newcounter{countdef}						% Zähler für Definitionen
\numberwithin{countdef}{section}			% Section in Definition-Nummerierung einbeziehen
\newcounter{countlemma}						% Zähler für Lemmata
\numberwithin{countlemma}{section}			% Section in Lemma-Nummerierung einbeziehen

% ab hier können neue Umgebungen für Sätze definiert und benannt werden!
% Syntax: \newtheorem{interner Umgebungsname}[Zähler]{Name im Dokument}
% Möchte man Sätze und Definitionen getrennt durchnummerieren, benötigt man weitere Zähler.
\newtheorem{satz}[countsatz]{Satz}
\newtheorem{bsp}[countbsp]{Beispiel}
\newtheorem{falko}[countfalko]{F\hspace{-0.25em}}
\newtheorem{defn}[countdef]{Definition}
\newtheorem{lemma}[countlemma]{Lemma}
% % % % % % % % % % % % % % % % % % % % % % %
\newcommand{\PP}{\mathbb{P}}
\DeclareMathOperator{\li}{li}
\newcommand{\kon}{\equiv}

% % % % % % % % % % % % % % % % % % % % % % %
\begin{document}
\begin{titlepage}
	% Nach einer Vorlage von http://www.LaTeXTemplates.com
	\newcommand{\HRule}{\rule{\linewidth}{0.8mm}} % Defines a new command for the horizontal lines, change thickness here

	\center % Center everything on the page
 
	\begin{minipage}{0.4\textwidth}
	\begin{flushleft}
	\includegraphics[height=1.5cm,keepaspectratio]{img/wwulogo.pdf}\\[1cm]
	\end{flushleft}
	\end{minipage}
	\hfill
	\begin{minipage}{0.4\textwidth}
	\begin{flushright}
	\vspace*{0.3cm}
	\includegraphics[height=1.2cm,keepaspectratio]{img/fb10logo.pdf} \
	\end{flushright}
	\end{minipage}

	\vspace{2cm}
	


	\HRule \\[0.8cm]
	{ \huge \sffamily\bfseries $\mathbf{CAT}\boldsymbol{(0)}$ kubische Komplexe}\\[0.4cm] % Title of your document
	\HRule \\[1cm]
 
	{\LARGE gelesen von} \\[.7cm]
	\textsc{\LARGE \textbf{Dr. Olga Varghese}}\\[.7cm]
	{\LARGE im Wintersemester 2015/2016}\\[2cm]

	\begin{tikzpicture}[scale=2]
		\draw[very thick,dashed] (1,1) -- (1.45,1.45) -- (2.45,1.45) -- (2,1);
		\draw[very thick,dashed] (2.45,1.45) -- (2.45,1);
		\draw[very thick,dashed] (1.45,1.45) -- (1.45,.45) -- (1,0);
		\draw[very thick,dashed] (1.45,.45) -- (2.45,.45);
		\draw[very thick,dashed] (3,1) -- (3.45,1.45) -- (4.45,1.45);
		\draw[very thick,dashed] (3.45,1.45) -- (3.45,2.45);
				
		\draw[very thick] (0,1) -- (2,1) -- (2,0) -- (2.45,.45) -- (2.45,1);
		\draw[very thick] (2,0) -- (1,0) -- (1,2) -- (4,2) -- (4,1) -- (3,1) -- (3,3) -- (2,3) -- (2,0);
		\draw[very thick] (2,1) -- (3,1) -- (3,0) -- (5,0);
		\draw[very thick] (4,0) -- (4,1) -- (4.45,1.45) -- (5.45,1.45) -- (5.45,2.45) -- (3.45,2.45) -- (3,2);
		\draw[very thick] (4,2) -- (4.45,2.45) -- (4.45,1.45);
		\draw[very thick] (4.45,2.45) -- (4.45,3.45);

		\draw[ultra thick,color=red,dashed] (3,1.5) -- (3.45,1.95) -- (4.45,1.95);		
		\draw[ultra thick,color=red] (1,1.5) -- (4,1.5) -- (4.45,1.95) -- (5.45,1.95);
	\end{tikzpicture}

	\vfill
	{\Large Vorlesungsmitschrift von Phil Steinhorst} \\[.5cm]
	{\large Stand: \today}
	
\end{titlepage}							% Titelseite ausgeben
\section*{Vorwort}
\label{sec:preface}
	Der vorliegende Text ist eine inhaltliche Aufbereitung zur Vorlesung Elliptische Kurven und Kryptographie, gelesen von PD Dr. Karin Halupczok an der WWU Münster im Sommersemester 2015. Der Inhalt entspricht weitestgehend dem handschriftlichen Skript, welches auf der Vorlesungswebsite bereitsgestellt wird. Dieses Werk ist daher keine Eigenleistung des Autors und wird nicht von der Dozentin der Veranstaltung korrekturgelesen. Für die Korrektheit des Inhalts wird keinerlei Garantie übernommen. Bemerkungen, Korrekturen und Ergänzungen kann man folgenderweise loswerden:
	\begin{itemize}
		\item persönlich durch Überreichen von Notizen oder per E-Mail
		\item durch Abändern der entsprechenden TeX-Dateien und Versand per E-Mail an mich
		\item direktes Mitarbeiten via GitHub. Dieses Skript befindet sich im \texttt{latex-wwu}-Repository von Jannes Bantje:
		\begin{center}
			\url{https://github.com/JaMeZ-B/latex-wwu}
		\end{center}
	\end{itemize}

\subsection*{Literatur}
\label{sub:lit}
\begin{itemize}
	\item Blake, Seroussi, Smart: Elliptic curves in cryptography
	\item Menezes, van Oorschot, Vanstone: Handbook of applied cryptography
	\item Silverman: The arithmetic of elliptic curves
	\item Silverman: A friendly introduction to number theory, chap. 40-45
	\item Washington: Elliptic curves, number theory and cryptography
	\item Werner: Elliptische Kurven in der Kryptographie
\end{itemize}

\subsection*{Kommentar der Dozentin}
In der Vorlesung beschäftigen wir uns mit den arithmetischen und geometrischen Eigenschaften elliptischer Kurven sowie deren Anwendungen in der Kryptographie. Dabei werden wir auch einen Vergleich mit Anwendungen der elementaren Zahlentheorie in der Kryptographie ziehen. Wir verfolgen eine elementare Herangehensweise, d.h. Kenntnisse der algebraischen Geometrie und der Funktionen- oder Zahlentheorie werden nicht benötigt. Es genügen die Vorkenntnisse aus den Grundvorlesungen.

\subsection*{Vorlesungswebsite}
\label{sub:link}
Das handgeschriebene Skript sowie weiteres Material findet man unter folgendem Link:
\begin{center}
	\url{\homepage}
\end{center}

\subsection*{Titelbild}
\label{sub:titlepic}
Das fehlt noch. Über Ideen und Anregungen freue ich mich sehr!

\vfill
\begin{flushright}
	Phil Steinhorst \\
	p.st@wwu.de
\end{flushright}
\newpage
\tableofcontents
\newpage											% Inhaltsverzeichnis ausgeben
\chapter{CAT(0)-Räume} % (fold)
\label{cha:1}

\begin{definition}[metrischer Raum]
	Sei $X$ eine Menge. Eine Abbildung $d \colon X \times X \rightarrow \RR_{\geq 0}$ heißt \Index{Metrik}, wenn für alle $x,y,z \in X$ gilt:
	\begin{enumerate}[(i)]
		\item $d(x,y) = 0 \Leftrightarrow x=y$
		\item $d(x,y) = d(y,x)$
		\item $d(x,z) \leq d(x,y) + d(y,z)$
	\end{enumerate}
	Das Paar $(X,d)$ heißt dann \Index{metrischer Raum}.
\end{definition}

\cleardoubleoddemptypage
\section{Der euklidische Algorithmus}
\label{sec:para2}
	Sei $R$ kommutativer Ring mit $1 \neq 0$. Für beliebiges $a \in R$ betrachte man die Menge der Vielfachen von $a \in R$, also
	\[ Ra := {xa : x \in R} = {b \in R : a|b} \]
	Die Teilmenge $I = Ra$ hat folgende Eigenschaften:
	\begin{enumerate}[(i)]
		\item $0 \in I$
		\item $b_1,b_2 \in I \Rightarrow b_1+b_2 \in I$
		\item $c \in R, b \in I \Rightarrow cb \in I$
	\end{enumerate}
	
\begin{defn}[Ideal, Hauptideal] \label{def_2.1}
	Eine Teilmenge $I$ von $R$ heißt ein \Index{Ideal} in $R$, falls die Eigenschaften (i), (ii), (iii) erfüllt sind. $I$ heißt \Index{Hauptideal}, wenn es ein $a \in R$ gibt mit $I = Ra$. Wir verwenden die Bezeichnung
	\[ (a) := Ra \]
	und nennen $(a)$ das von $a \in R$ erzeugte Hauptideal.
\end{defn}

\minisec{Bemerkung}
	\begin{enumerate}[(1)]
		\item $(b) \subseteq (a) \Leftrightarrow a | b$
		\item $a \assoz b \Leftrightarrow (a) = (b)$
		\item $c$ ist gemeinsames Vielfaches von $a_1,\dots,a_n \Leftrightarrow (c) \subseteq (a_1) \cap \dots \cap (a_n)$
		\item $m$ ist ein kgV von $a_1,\dots,a_n \Leftrightarrow (a_1) \cap \dots \cap (a_n) = (m)$
		\item $d$ ist ein gemeinsamer Teiler von $a_1,\dots,a_n \Leftrightarrow (a_i) \subseteq (d)$ für $1 \leq i \leq n$
		\item $d$ ist ein gemeinsamer Teiler von $a_1,\dots,a_n \Leftrightarrow Ra_1 + Ra_2 + \dots + Ra_n \subseteq (d)$
		\item $d$ ist ein ggT von $a_1,\dots,a_n \Leftrightarrow (d)$ ist das kleinste Hauptideal mit $Ra_1 + \dots Ra_n \subseteq (d)$.
	\end{enumerate}
	Ein ggT lässt sich also idealtheoretisch nicht so einfach charakterisieren wie oben ein kgV durch (4). Am schönsten wäre es, wenn $Ra_1 + \dots + Ra_n$ ein Hauptideal wäre, dann würde (7) übergehen in:
	\[ d \text{ ist ein ggT von } a_1,\dots,a_n \Leftrightarrow Ra_1+Ra_2+ \dots + Ra_n = (d) \]
	
\begin{defn}[Hauptidealring] \label{def_2.2}
	Ein Integritätsring $R$ heißt ein \Index{Hauptidealring}, wenn jedes Ideal $I$ von $R$ ein Hauptideal ist.
\end{defn}

\minisec{Bezeichnung}
	Für Elemente $a_1,\dots,a_n$ in einem beliebigen kommutativen Ring $R$ mit $1 \neq 0$ setze
	\[ (a_1,\dots,a_n) := Ra_1 + \dots + Ra_n \]
	Man nennt $(a_1,\dots,a_n)$ das von $a_1,\dots,a_n$ erzeugte Ideal in $R$.
\newpage
\section{Partielle Differentialgleichungen erster Ordnung}
\label{sec:para3}

% % % % % % 22. Apr.
\minisec{Motivation: Linearer Fall}
	Betrachte \marginnote{22. Apr}
	\[ a(x) \cdot \nabla u(x) + a_0(x) = 0 \] mit Randdaten auf glatter $(n-1)$-dimensionaler Hyperfläche $\Gamma \subset \RR^n$, $u(x) = g(x)$ auf $\Gamma$. \\
	Die Richtungsableitung von $u$ in Richtung $a$ ist $a(x) \cdot \nabla u(x) = -a_0 (x)$, d.h. wir wissen, wie sich $u$ entlang charakteristischer Kurven $s \mapsto x(s)$ ändert mit
	\begin{equation}
	\begin{aligned}
		\dot{x}(s) = a(x(s)) \label{eq_2204_1}
	\end{aligned}
	\end{equation} \\
	Sei $\Gamma$ parametrisiert durch $x_0 \colon \RR^{n-1} \rightarrow \Gamma$. Zu $\widehat{x} \in \RR^n$ können wir $s \in \RR, y \in \RR^{n-1}$ finden, sodass $\widehat{x} = x(s)$ für eine Lösung $x(s)$ von \eqref{eq_2204_1} mit $x(0) = x_0(y)$, d.h. wir reparametrisieren $\RR^n$ durch eine Funktion $\widehat{x}\colon \RR^{n-1} \times \RR \rightarrow \RR, (y,s) \mapsto \widehat{x}(y,s)$. Dann gilt:
	\[ u(\widehat{x}(y,s)) = \underbrace{u(\widehat{x}(y,0))}_{g(\widehat{x}(y,0))} + \int\limits_{0}^{s} \underbrace{\frac{\der u(\widehat{x}(y,t))}{\der \tau}}_{=-a_0(\widehat{x}(y,\tau))} d\tau; \quad u(x) = (u \circ \widehat{x}) \circ \widehat{x}^{-1} \]
	
\begin{bem} \label{bem_2}
	Damit $u$ auf $\RR^n$ definiert ist, muss $\widehat{x}^{-1}$ definiert sein.\marginnote{[2]} Für lokale Invertierbarkeit ist hinreichend:
	\[ \det(D \widehat{x}) \neq 0 \]
	Auf $\Gamma$ ist $\det(D \widehat{x}(y,0) = \det( \underbrace{D x_0(y)}_{\in R^{n\times m -1}}\vert a(x_0(y)))$, d.h. wir brauchen $a(x_0(y)) \notin \operatorname{span}(Dx_0(y))$. Die Spalten von $Dx_0(y)$ spannen den Tangentialraum von $\Gamma$ an $x_0(y)$ auf. Damit dies ungleich 0 ist, muss $a \notin \operatorname{span}(Dx_0)$, d.h. $a \cdot \nu \neq 0$ mit $\nu$ als Normale auf $\Gamma$, d.h. $\Gamma$ ist nichtcharakteristisch.
\end{bem}
	
\begin{bsp} \label{bsp_3}
	Betrachte \marginnote{[3]}
	\[ \begin{cases}
		u_{x_1} + u_{x_2} = 1 \text{ auf } \RR^2 \\
		u = 0 \text{ auf } \Gamma = \{x_1 + x_2 = 0 \}
	\end{cases} \]
	Charakteristische Kurven: $\dot{x} = \begin{pmatrix} 1 \\ 1	\end{pmatrix} \Rightarrow x(s) = \begin{pmatrix} c_1 + s \\ c_2 + s \end{pmatrix}$. \\
	Parametrisiere $\Gamma$:
	\[ x_0(y) := \begin{pmatrix} y \\ -y \end{pmatrix} = x_0(y) \]
	$\Rightarrow \widehat{x}(y,s) = x(s)$ für die Kurve $x(s)$ mit $x(0) = x_0(y) \Rightarrow \widehat{x}(y,s) = \begin{pmatrix} y+s \\ -y+s \end{pmatrix}$ \\
	$\Rightarrow \widehat{x}^{-1}(x) = (\frac{x_1-x_2}{2},\frac{x_1+x_2}{2})$ \\
	$\Rightarrow u(\widehat{x}(y,s)) = 0 + \int\limits_{0}^{s} 1 d\tau = s$ \\
	$\Rightarrow u(x) = \underbrace{(u \circ \widehat{x})}_s \circ \widehat{x}^{-1}(x) = \frac{x_1+x_2}{2}$
\end{bsp}
	
\minisec{Nichtlinearer Fall}
	Betrachte
	\begin{equation}
		0 = F(\nabla u(x),u(x),x), \quad u=g \text{ auf } \Gamma, \quad F,g,\Gamma \ C^2\text{-glatt} \label{eq_2204_2}
	\end{equation}
	Idee: Information fließt entlang charakteristischer Kurven \\
	Ziel: Finde charakteristische Kurven und wandele die PDGL um in ein System von gewöhnlichen DGL entlang der Kurven. \\
	Notation: \begin{itemize}
	\item Für eine Kurve $x(s)$ schreibe $z(s) := u(x(s))$ und $p(s) = \nabla u(x(s))$, d.h.
	\[ 0 = F(p(s), z(s), x(s)) \]
	\item $\dot{} = \frac{\der}{\der s}$
	\item $F_p,F_z,F_x$ sind partielle Ableitungen von $F$ nach erstem, zweitem und drittem Argument.
	\end{itemize}
	Suche GDGL für $x,p,z$.
	
\begin{thm}[Charpits Gleichungen]
\label{thm_charpits} \label{thm_4}
	Sei $u \in C^2(\Omega)$ eine Lösung von \eqref{eq_2204_2}.\marginnote{[4]} Entlang von Kurven $x(s)$ in $\Omega$ mit
	\begin{equation}
	\begin{aligned}
		\dot{x} = F_p \label{eq_2204_3}
	\end{aligned}
	\end{equation}
	gilt
	\begin{align}
		&\dot{z} = p \cdot F_p \label{eq_2204_4} \\
		&\dot{p} = -pF_z - F_x \label{eq_2204_5}
	\end{align}
	Diese GDGLs heißen \Index{Charpits} oder \Index{charakteristische Gleichungen}, Kurven $x(s)$ oder $(x(s),z(s),p(s))$ heißen \Index{charakteristische Kurven}.
\end{thm}
	
\minisec{Beweis}
	Betrachte beliebige Kurve $x(s)$. Dann gilt:
	\[ \begin{array}{c}
		\dot{z}(s) = \nabla u(x(s)) \cdot \dot{x}(s) = p(s) \cdot \dot{x}(s) \\
		\dot{p}(s) = D^2u(x(s)) \dot{x}(s)
	\end{array} \]
	Um $D^2u$ loszuwerden, differenziere \eqref{eq_2204_2} nach $x$:
	\[ 0 = \frac{d}{dx} F(\nabla u(x),u(x),x) = D^2 u F_p + F_z \overbrace{\nabla u}^p + F_x \]
	Wähle die Kurve $x(s)$, sodass $\dot{x}(s) = F_p(p(s),z(s),x(s))$, so erhalten wir
	\[ \dot{p}(s) = -F_z p - F_x \qquad \dot{z} = p \cdot F_p \]. \qed

\mbox{} \\
\eqref{eq_2204_3}, \eqref{eq_2204_4} und \eqref{eq_2204_5} können mit Anfangswerten $x(0),z(0),p(0)$ gelöst werden. Sei $\Gamma$ parametrisiert durch $x_0(y), y \in \RR^{n-1}$. Werte für $x(0)$ und $z(0)$ sind auf $\Gamma$ gegeben durch $x(0) = x_0(y) \in \Gamma$ und $z(0) = g(x_0(y))$. $p(0)$ erhalten wir durch Lösen von
\begin{equation}
	\left. \begin{cases}
	0 = F(p(0),z(0),x(0)) & \text{PDGL} \\
	0 = \nabla x_0 (y) p(0) - \nabla_y g(x_0(y)) & \text{Abl. von } u(x_0(y)) = g(x_0(y)) \text{ nach } y
	\end{cases} \right\}  \label{eq_2204_7} \end{equation}
wobei
	\[ \nabla x_0 (y) = \begin{pmatrix}
	\der_{y_1} {x_0}_1 & \cdots & \der_{y_1} {x_0}_n \\ 
	\vdots & \ddots & \vdots \\ 
	\der_{y_{n-1}} {x_0}_1 & \cdots & \der_{y_{n-1}} {x_0}_n
	\end{pmatrix} \]
Die Lösung von \eqref{eq_2204_3}, \eqref{eq_2204_4} und \eqref{eq_2204_5} für diese Anfangswerte liefert eine Parametrisierung $\widehat{x}(y,s)$ des $\RR^n$ mit $\widehat{x}(y,s) = x(s)$ für die Kurve $x(s)$ mit $x(0) = x_0(y)$ sowie Funktionen $z(y,s),p(y,s)$.

\begin{defn}[nichtcharakteristische Hyperfläche] \label{def_nichtchar} \label{def_5}
	Eine glatte $(n-1)$-dimensionale Hyperfläche\marginnote{[5]} $\Gamma \subset \RR^n$ heißt \Index{nichtcharakteristisch} zu \eqref{eq_2204_2}, wenn für die Einheitsnormale $\nu$ an $\Gamma$ gilt:
	\[ F_p(p,z,x) \cdot \nu(x) \neq 0 \text{ für alle } x \in \Gamma, z \in \RR, p \in \RR^n \]
\end{defn}
	
\begin{lemma} \label{lemma_6}
	Wenn gilt: \marginnote{[6]} \begin{itemize}
	\item $\Gamma$ ist nichtcharakteristisch
	\item Für ein gegebenes $y_0 \in \RR^{n-1}$ existiert eine Lösung von \eqref{eq_2204_7}
	\end{itemize}
	dann: \begin{itemize}
		\item kann \eqref{eq_2204_7} auch in einer Umgebung $U$ von $y_0$ nach $p(0)$ gelöst werden
		\item existieren ein offenes Intervall $I = (-\delta,\delta)$ und eine Umgebung $V$ von $x_0(y_0)$, sodass für alle $x \in V$ ein $y \in U$ und $s \in I$ existiert mit $x = \widehat{x}(y,s)$. $\widehat{x}^{-1}$ ist $C^2$.
	\end{itemize}
\end{lemma}
	
\minisec{Beweis}
	\begin{itemize}
		\item Existenz von $p(0)$ folgt direkt aus dem Satz über implizite Funktionen. (Beachte:
		\[ \frac{\der}{\der p(0)} \begin{pmatrix}
		F(p(0),z(0),x(0)) \\ \nabla x_0 (y) p(0) - \nabla_y g(x_0(y)) \end{pmatrix} = \begin{pmatrix}
		F_p \\ \nabla x_0 (y) \end{pmatrix} \]
		hat vollen Rang.)
		\item Existenz und Glattheit von $\widehat{x}^{-1}$ folgt aus dem Satz über inverse Funktionen, da $\det (y_0,0) = \det \begin{pmatrix}	\nabla x_0 \\ F_p \end{pmatrix} \neq 0$ \qed
	\end{itemize}
	
\begin{bem} \label{bem_8}
	Auch hier bedeutet nichtcharakteristisch,\marginnote{[7]} dass auf $\Gamma$ nach den partiellen Ableitungen von $u$ gelöst werden kann.
\end{bem}
	
\begin{thm}[Lokale Existenz einer glatten Lösung]
\label{thm_lok_existenz} \label{thm_8}
	Seien \marginnote{[8]} \begin{itemize}
		\item $\Gamma, g, F$ glatt, also $C^2$
		\item $\Gamma$ nichtcharakteristisch und
		\item es gebe ein $y_0 \in \RR^{n-1}$ und ein $p(0) \in \RR^n$, das \eqref{eq_2204_7} löst.
	\end{itemize}
	Dann existiert auf der Umgebung $V$ von $x_0(y_0)$ aus Lemma \ref{lemma_6} eine $C^2$-Lösung $u$ von \eqref{eq_2204_2} und \eqref{eq_2204_3}. Es ist $u(x) = z(y,s)$ und $\nabla u(x) = p(y,s)$ für $(y,s) = \widehat{x}^{-1} (x)$.
\end{thm}

% % % % % % % % % % % % % 25.04.
\minisec{Beweis}
	\begin{enumerate}[1.]
		\item Nach Lemma \eqref{lemma_6} \marginnote{25. Apr} existiert $\widehat{x}^{-1}$ auf $V$ sowie $z(y,0) = g(x_0(y)),p(y,0)$ mit $F(p(y,0),z(y,0),\widehat{x}(y,0)) = 0$ und $\nabla x_0 (y) p(0) - \nabla_y g(x_0(y)) = 0$. Löse Charpits Gleichungen mit diesen Anfangswerten $\Rightarrow z(y,s),p(y,s)$.
		\item Zeige: $f(y,s) := F(p(y,s),z(y,s),\widehat{x}(y,s)) = 0$. \\
			$f(y,0) = 0, \dot{f}(y,s) = F_p \cdot \dot{p} + F_z \cdot \dot{z} + F_x \dot{\widehat{x}} = F_p(-pF_z - F_x) + F_zpF_p + F_x \cdot F_p = 0$
		\item Zeige $p(\widehat{x}^{-1}(x)) = \nabla u(x)$ für $u(x) = z(\widehat{x}^{-1}(x))$. \marginnote{$z \equiv u \circ \widehat{x}^{-1}$ \\ $z_s = \nabla u \cdot \widehat{x}_s = p \cdot \widehat{x}_s$ \\ $z_y = \nabla u \widehat{x}_y = p \cdot \widehat{x}_y$} \begin{itemize}
				\item $\cdot{z}(y,s) = p(y,s) \cdot \dot{\widehat{x}}(y,s)$
				\item $z_y(y,s) = p^T(y,s) \widehat{x}_y(y,s)$, denn $r(s) := z_y(y,s) - p^T(y,s) \widehat{x}_y (y,s)$ erfüllt
				\begin{equation}
				\begin{aligned}
					\dot{r}(s) &= \underbrace{\dot{z_y}(y,s)}_{p^T_y \dot{\widehat{x}} + p^T \dot{\widehat{x}}_y} - \dot{p}^T(y,s) \widehat{x}_y(y,s) - p^T(y,s) \dot{\widehat{x}}_y(y,s) = p_y^T F_p - (-pF_z - F_x)^T \widehat{x}_y \\ \notag
					&= \frac{\der}{\der y} \underbrace{F(p(y,s),z(y,s),\widehat{x}(y,s))}_{=0} + \underbrace{p^T F_z \widehat{x}_y - F_z z_y}_{-F_z r}
				\end{aligned}
				\end{equation} 
				mit Anfangswert $r(0) = \underbrace{z_y(y,0)}_{\nabla_y g(x_0(y))} - p^T(y,0) \underbrace{\widehat{x}_y(y,0)}_{\nabla_y x_0(y)} = 0$. $\Rightarrow r \equiv 0$.
				\item $u_x = z_y y_x + z_s s_x = p^T \widehat{x}_y y_x + p^T \widehat{x}_s s_x = p^T \underbrace{(\widehat{x}_y y_x + \widehat{x}_s s_y)}_{\widehat{x}_x} = p^T, p = \nabla u$. \qed
			\end{itemize}
	\end{enumerate}

\begin{bsp} \label{bsp_9}
	Löse das Cauchy-Problem \marginnote{[9]}
	\[ \begin{cases}
		u_t + uu_x = 1 & \text{ auf } (0,\infty) \times \RR \\
		u = x & \text{ auf } \Gamma = \{t = 0\} \end{cases} \]
	Charakteristische Kurven sind gegeben durch $\frac{\der}{\der s} \begin{pmatrix} (t,x) \\ z \\ 	p \end{pmatrix} = \begin{pmatrix} (1,z) \\ P-1+zp_2 \\ -p_2 p \end{pmatrix}$, also
	\[ \begin{pmatrix} (t,x) \\ z \\ 	p \end{pmatrix} = \begin{pmatrix} (t(0)+s,x(0) + z(0)s + \frac{p_1(0)+z(0)p_2(0)}{2} s^2 \\ p_1(0)s + z(0)(1+p_2(0)s) \\ \frac{p}{1+p_2(0)s} \end{pmatrix}.\]
	Parametrisiere $\Gamma$ durch $x_0(y) = (0,y)$. Das liefert Anfangswerte $(t(0),x(0))=x_0(y)=(0,y)$ und $z(0) = y$. Lösen von \eqref{eq_2204_7} liefert $p(0) = \begin{pmatrix} 1-y \\ 1 \end{pmatrix}$. Insgesamt ergibt sich:
	\[ \Rightarrow \begin{pmatrix} (t,x) \\ z \\ p \end{pmatrix} = \begin{pmatrix} s,y+ys+\frac{s^2}{2} \\ y+s \\ \frac{1}{1+s} \begin{pmatrix} 1-y \\ 1 \end{pmatrix} \end{pmatrix} \Rightarrow (s,y) = \left( t,\frac{x-\frac{t^2}{2}}{1+t}\right) \Rightarrow u(t,x) = z(y,s) = y+s = t + \frac{x-\frac{t^2}{2}}{1+t} \]
\end{bsp}
	
\begin{bem} \label{bem_10}
	Für quasilineare PDGL der Form \marginnote{[10]}
	\[ b(u(x),x) \cdot \nabla u(x) + c(u(x),x) = 0 \]
	können die charakteristischen Gleichungen \eqref{eq_2204_3}, \eqref{eq_2204_4} und \eqref{eq_2204_5} vereinfacht werden: $b(z(s),x(s)) \cdot p(s) + c(z(s),x(s)) = 0$ impliziert:
	\begin{align}
		\dot{x}(s) = b(z(s),x(s))	\label{eq_2504_8} \\
		\dot{z}(s) = -c(z(s),x(s))	\label{eq_2504_9} \\
		\dot{p}(s) = -(b_z \cdot p + c_z)\cdot p - b_x p -c_x	\label{eq_2504_10}
	\end{align}
	Die ersten beiden Gleichungen legen bereits $x$ und $z$ fest, sodass \eqref{eq_2504_10} ignoriert werden kann.
\end{bem}
\newpage
\section{Sichere Kryptographie mit elliptischen Kurven}
\label{sec:para4}
\nextlecture
\subsection{Bekannte Angriffe auf das DL-Problem: Überblick}
\label{sub:4.1}
\begin{bem}
	Die Sicherheit des ElGamal- und ECDSA-Verfahrens beruht hier auf der Schwierigkeit des DL-Problems auf elliptischen Kurven. \marginnote{[19]}
	Allerdings gibt es bestimmte Arten elliptischer Kurven, bei denen das DL-Problem algorithmisch schnell lösbar ist, sodass sich diese Kurven als kryptographisch schwach bzw. ungeeignet erweisen.
	Auch die Wahl eines Punktes $P$ mit großer Ordnung ist wichtig.
\end{bem}

\subsubsection{BSGS und Silver-Pohlig-Hellman}
\label{subsub:4.1.1}
	Diese beiden Methoden eignen sich zur Lösung des DL-PRoblems in einer beliebigen abelschen Gruppe $G$.
	
\begin{defn}[DL-Problem in $G$]
	Gegeben sei $P \in G$ mit $\ord(P) = n \in \NN$, sowie $Q \in \sprod{P}$. Gesucht ist $k \in \{0,\dots,n-1\}$ mit $kP = Q$.
\end{defn}

\begin{defn}[BSGS -- Baby Steps Giant Steps]
	Dieses Verfahren kommt in Frage, wenn $P$ kleine Ordnung $n$ hat.
	Dann kann das DL-Problem wie folgt gelöst werden; der Algorithmus hat einen Zeit- und Platzbedarf der Größenordnung $\oh(\sqrt{n})$: \index{BSGS (Baby Steps Giant Steps)}
\end{defn}

\begin{bem}
	Vorüberlegung: Sei $m = \ceil*{\sqrt{n}} = \min \{ l \in \NN : l \geq \sqrt{n}\}$, schreibe $k = qm + r, r \in \{0,1,\dots,m-1\}$ (Division mit Rest). \\
	Ziel: Bestimme $q, r$. \\
	Da $Q = kP = qmP + rP$, folgt $\underbrace{Q - rP}_{\text{"Baby step"}} = \underbrace{qmP}_{\text{"Giant step"}}$.
\end{bem}

\begin{bem}
	Idee: Berechne alle möglichen Werte von "Baby step" und nach und nach die möglichen Werte von "Giant step".
	Trifft man auf eine Übereinstimmung, sind $r$ und $m$ gefunden.
\end{bem}

\begin{anw}[Schritt 1]
	Berechne die Liste der "Baby steps" $B = \{(Q-rP,r) : 0 \leq r \leq m \}$.
\end{anw}

\begin{anw}[Schritt 2]
	\begin{itemize}
		\item Ist für eines der $r$ die Gleichung $Q - rP = \oh$ erfüllt, ist $k = r$.
		\item Sonst teste für den ersten "Giant step" $R = mp$, ob $R$ in $B$ schon vorkommt.
		Falls ja: $k = m + r$.
		\item Teste so alle "Giant steps" $2R, 3R, 4R, \dots, (m-1)R$, ob diese in $B$ vorkommt.
		Wenn ja, gibt die zweite Komponente $r$ mit $k = qm+r$.
	\end{itemize}
\end{anw}

\begin{defn}[Silver-Pohlig-Hellman-Verfahren]
	Dieses Verfahren löst das DL-Problem in einer abelschen Gruppe $G$, wenn die Ordnung $n = \ord(P)$ aus nur kleinen Primfaktoren $p_i$ zusammengesetzt ist, d.h. \Index{glatt} ist. \index{Silver-Pohlig-Hellman}
\end{defn}

\begin{defn}[$B$-glatt]
	Sei $B \in \RR_{>0}$.
	Dann heißt $m \in \NN$ \bet{$B$-glatt}, falls für alle $p \mid n$ gilt: $p \leq B$.
	\index{B-glatt@$B$-glatt}
\end{defn}

\begin{bem}
	Die Bestimmung von $k$ in $\sprod{P}$ wird auf Untergruppen von $\sprod{P}$ der Ordnungen $p_i \mid n$ zurückgeführt.
	Sei $\ord(P) = n = \prod_{i=1}^{t} p_i^{\lambda_i}$ mit $p_1, \dots, p_t$ paarweise verschieden und prim, $\lambda_i \in \NN$.
	Der Algorithmus hat dann eine Laufzeit von $\oh \enbrace*{\sum_{i=1}^{t} (\lambda_i(\log n + \sqrt{p_i}))}$.
\end{bem}

\begin{bem}
	Vorüberlegung:
	\begin{itemize}
		\item Zur Bestimmung von $k$ mit $kP = Q \in \sprod{P}$ berechnen wir alle Restklassen $k \modu p_1^{\lambda_1}, k \modu p_2^{\lambda_2}, \dots, k \modu p_t^{\lambda_t}$.
		Denn laut chinesischem Restsatz ist $\ZZ/n\ZZ \simeq \ZZ/p_1^{\lambda_1} \ZZ \times \dots \times \ZZ / p_t^{\lambda_t} \ZZ$, sodass damit dann auch die Restklasse von $k \modu n$ bestimmt werden kann.
		\item Betrachte daher jedes $p = p_i, \lambda = \lambda_i$ mit $1 \leq  i \leq t$. \\
		Gesucht: $z \in \{0, \dots, p^\lambda - 1\}$ mit $z \kon k \modu p^\lambda$. \\
		Schreibe $z = z_0 + z_1 p + \dots + z_{\lambda-1} p^{\lambda-1}$, die $z_i \in \{0, \dots, p-1\}$, in der $p$-adischen Entwicklung; bestimme die $z_0, \dots, z_{\lambda-1}$.
	\end{itemize}
\end{bem}

\begin{anw}[Schritt 1]
	Sei $R := \frac{n}{p} \cdot P$, dann ist $\frac{n}{p}Q = \frac{n}{p}kP = kR$ und $pR = \oh$.
	Also ist $kR = zR = z_0 R$, d.h. $z_0 R = \frac{n}{p} \cdot Q$.
	Somit muss man in der Untergruppe $\sprod{R}$ der (kleinen) Ordnung $p$ ein DL-Problem lösen, um $z_0$ zu bestimmen -- etwa mit BSGS.
\end{anw}

\begin{anw}[Schritt 2]
	Seien $z_0, \dots, z_{j-1}$ schon (rekursiv) bestimmt, wo $j \leq \lambda-1$ ist.
	Berechne dann $Q_j := \frac{n}{p^{j+1}} (Q - (z_0 + z_1p + \dots + z_{j-1} p^{j-1}) P)$. Da $nP = \oh$, ist $\frac{n}{p^{j+1}} \cdot p^\lambda P = \oh$; da $z \kon k \modu p^\lambda$, ist $k = z + sp^\lambda, s \in \ZZ$, also $\frac{n}{p^{j+1}} Q = \frac{n}{p^{j+1}} k P = \frac{n}{p^{j+1}} zP + \underbrace{\frac{n}{p^{j+1}} \cdot sp^\lambda P}_{=\oh} = \frac{n}{p^{j+1}} zP$ und somit $Q_j = \frac{n}{p^{j+1}} (z_j p^j + \dots + z_{\lambda - 1} p^{\lambda-1}) P = \frac{n}{P} z_j P = z_j R$. \\
	Zur Berechnung von $z_j$ ist wieder ein DL-Problem in der Untergruppe $\sprod{R}$ der Ordnung $p$ zu lösen -- etwa mit BSGS.
\end{anw}

\begin{bem}
	Ist die Gruppenordnung glatt, ist der Algorithmus also sehr schnell.
\end{bem}

\subsubsection{Pollard-$\rho$ und Pollard-$\lambda$}
\label{subsub:4.1.2}
\begin{bem}
	Der \bet{Pollard-$\rho$-Algorithmus} ist von der Laufzeit her vergleichbar mit BSGS, ist aber speicherplatztechnisch günstiger und lässt sich gut parallelisieren.
	Mit $m$ Prozessoren wird der Algorithmus so um den Faktor $m$ schneller. \index{Pollard-$\rho$}
\end{bem}

\begin{bem}
	Der \bet{Pollard-$\lambda$-Algorithmus} ist ähnlich, im Allgemeinen aber eher langsamer als Pollard-$\rho$.
	Er liefert gute Ergebnisse, wenn der diskrete Logarithmus in einem hinreichend kleinen Intervall liegt.
	Auch Pollard-$\lambda$ ist gut parallelisierbar.
	(Die genauen Verfahren können in der Fachliteratur nachgeschlagen werden.) \index{Pollard-$\lambda$}
\end{bem}

\subsubsection{MOV und SSSA}
\label{subsub:4.1.3}
\begin{bem}
	Beim \Index{MOV-Verfahren} (Autoren: Menezes, Okamoto, Vanstone) wird das DL-Problem für eine elliptische Kurve $E(\FF_{p^r})$ auf das in der Gruppe $(\FF_{p^{rl}}^*,\cdot)$ für ein $l \geq 1$ zurückgeführt.
	Es ist also speziell nur für elliptische Kurven konstruiert, nicht für allgemeine abelsche Gruppen.
	Zeigt sich hier, dass $l \geq 1$ so wählbar ist, dass das DL-Problem in $(\FF_{p^{rl}}^*,\cdot)$ leicht, d.h. schnell, zu lösen ist, ist die elliptische Kurve kryptographisch ungeeignet, etwa wenn $n = \ord(P)$ Teiler von $p^{rl} - 1$ ist.
\end{bem}

\begin{bem}
	Generell lässt sich das DL-Problem in $(\FF_{p^{rl}}^*,\cdot)$ in subexponentieller Zeit schnell lösen (mit so genannten Indexkalkül-Methoden), sodass Kurven, für die das DL-Problem auf ein schnelles in einem $(\FF_{p^{rl}}^*,\cdot)$ zurückgeführt werden kann, als kryptographisch schwach bzw. ungeeignet angesehen werden.
	Das ist etwa bei supersingulären elliptischen Kurven der Fall, bei denen die Gruppenstruktur recht gut bekannt ist.
\end{bem}

\begin{defn}[supersingulär]
	Eine elliptische Kurve $E(\FF_{p^r})$ heißt \Index{supersingulär}, falls $ p =\Char(\FF_{p^r})$ die Spur des Frobenius teilt, d.h. $p \mid p^r +1 - \#E(\FF_{p^r})$.
\end{defn}

\begin{bem}
	\begin{itemize}
		\item Um zu testen, ob eine Kurve supersingulär und damit kryptographisch ungeeignet ist, muss die Gruppenordnung $\#E(\FF_{p^r})$ der elliptischen Kurve bestimmt werden -- typischerweise mit dem Schoof-Algorithmus.
		\item Der Begriff supersingulär hat nichts mit singulären Punkten zu tun:
		elliptische Kurven sind per Definition nicht-singulär.
	\end{itemize}
\end{bem}

\begin{bsp}
\label{bsp_19.21}
	Die Kurve $E(\FF_2)\colon y^2 + y = x^3 + x + 1$ ist supersingulär, da $E(\FF_2) = \{ \oh \}$.
\end{bsp}

\begin{bem}
	Supersingularität bleibt bei Übergang zu einem Erweiterungkörper erhalten:
	Ist $E(\FF_{p^r})$ supersingulär, dann auch $E(\FF_{p^{rl}})$ für alle $l \geq 1$. (ohne Beweis)
\end{bem}

\begin{satz}[Erstes Kriterium für Supersingularität]
	Sei $p \geq 3$, $E(\FF_p) \colon y^2 = x^3 + ax^2 + bx + c =: h(x)$ elliptische Kurve.
	Dann ist $E(\FF_{p})$ genau dann supersingulär, wenn der Koeffizient vor $T^{p-1}$ ind $h(T)^{\frac{p-1}{2}} \in \FF_p[T]$ gleich $0$ ist.
\end{satz}

\begin{satz}[Zweites Kriterium für Supersingularität]
	Sei $p = 2$, $E(\FF_{2^r}) \colon y^2 + a_1 xy + y = x^3 + a_2 x^2 + a_4 x + a_6$.
	Dann ist $E(\FF_{2^r})$ genau dann supersingulär, wenn $a_1 = 0$. (ohne Beweis)
\end{satz}

\begin{bsp}
	Siehe Beispiel~\ref{bsp_19.21}, und $E(\FF_p) \colon y^2 = x^3 + x$ ist für $p \kon 3 \modu 4$ supersingulär, denn: \\
	$(T^3 + T)^{(p-1)/2} = \sum_{j=1}^{(p-1)/2} \binom{(p-1)/2}{j} T^{3j} T^{(p-1)/2-j}$ mit $\frac{p-1}{2} + 2j = p-1 \Leftrightarrow 2j = \frac{p-1}{2}$, d.h. wenn $2 \mid \frac{p-1}{2} \Leftrightarrow p \kon 1 \modu 4$ \\
	$\Rightarrow$ Koeffizient vor $T^{p-1}$ ist $\binom{(p-1)/2}{(p-1)/4} \neq 0$ in $\FF_p$. Für $p \kon 3 \modu 4$ kommt $T^{p-1}$ nicht vor, also Koeffizient ist $0$.
\end{bsp}

\begin{bem}
	Der MOV-Algorithmus nutzt bei einer supersingulären Kurve $E(\FF_{p^r})$ aus, dass $t = p^r + 1 - \#E(\FF_{p^r})$ nur einen der Werte $t \in \{0, \pm \sqrt{p^r}, \pm \sqrt{2p^r}, \pm \sqrt{3p^r}, \pm 2 \sqrt{p^r}\}$ annehmen kann.
\end{bem}

\begin{bem}
	Beim \Index{SSSA-Verfahren} (Autoren: Satoh, Smart, Semaev, Araki) handelt es sich um einen schnellen Algorithmus zur Lösung des DL-Problems auf anormalen elliptischen Kurven, welche deswegen kryptographisch ungeeignet sind.
	Die Grundidee ist, die elliptische Kurve über $\FF_p$ als eine über $\QQ_p$ zu betrachten, dem Körper der $p$-adischen Zahlen, und die Logarithmen auf eine Division in $\ZZ_p$ zurückzuführen (was leicht ist).
\end{bem}

\begin{defn}[anormal]
	Eine elliptische Kurve $E(\FF_p)$ heißt \Index{anormal}, wenn $\#E(\FF_p) = p$ ist.
	Dies lässt sich wieder durch Bestimmung von $\#E(\FF_p)$ mit dem Schoof-Algorithmus leicht überprüfen.
	Der SSSA-Algorithmus kann auf Kurven über $\FF_{p^r}$ übertragen werden.
	Er hat polynomielle Laufzeit.
\end{defn}

\subsubsection{Fazit: geeignete elliptische Kurven und Vergleich mit anderen Public-Key-Verfahren}
\begin{bem}
	Eine elliptische Kurve $E(\FF_{p}): y^2 = x^3 + ax + b$ mit vorgegebener Bitzahl für $p$ ist leicht zu finden -- mit Zufallszahlengenerator und Primzahltest, was auch für große Zahlen mit mehreren hundert Dezimalstellen schnell machtbar ist; dafür kennt man ganz gute Algorithmen.
\end{bem}

\begin{bem}
	Man wählt so lange die Parameter $p,a,b$ neu, bis die Diskriminante $4a^3+27b^2$ nicht durch $p$ teilbar ist und somit eine elliptische Kurve vorliegt.
	Ziemlich sicher liegt dann eine kryptographisch geeignete Kurve vor.
	Das testet man nach Berechnen der Gruppenordnung $\#E(\FF_p)$ mit dem Schoof-Algorithmus:
\end{bem}

\begin{bem}
	\begin{itemize}
		\item Ist $\#E(\FF_p)$ glatt, d.h. hat $\#E(\FF_p)$ nur kleine Primteiler, ist die Kurve ungeeignet wegen Silver-Pohlig-Hellman.
		\item Ist $\#E(\FF_p) = p+1$, d.h. die Kurve ist supersingulär, dann ist die Kurve ungeeignet (MOV).
		\item Ist $\#E(\FF_p) = p$, d.h. die Kurve anormal, ist die Kurve ungeeignet (SSSA).
	\end{itemize}
	Ob die Kurve supersingulär oder anormal ist, kann man meist leicht daran erkennen durch Wahl von Punkten $P \in E(\FF_p)$ und dem Test, ob $(p+1)P = \oh$ bzw. $pP = \oh$ gilt.
\end{bem}

\begin{bem}
	Die Wahl eines Punktes $P$ mit nicht zu kleiner Ordnung $n$ muss dann gewährleistet werden.
	Speziell darf $n$ kein Teiler von $p^{rl} - 1$ sein, wenn das DL-Problem in $(\FF_{p^{rl}}^*,\cdot)$ leicht zu lösen ist, und $n$ darf auch kein Vielfaches von $p$ sein (wegen SSSA).
	Auch sollte $n$ nicht glatt sein; man wählt in der Praxis meist Punkte $P$, für die $n = \ord(P)$ eine hinreichend große Primzahl ist; für sie sollte etwa $n > 2^{160}$ gelten.
\end{bem}

\begin{bem}
	Die für allgemeine elliptische Kurven, die in diesem sinne als kryptographisch sicher gelten, bekannten Implementationen des DL-Problems sind alle von exponentieller Komplexität.
	Ein Kryptographieverfahren wie ElGamal bzw. DSA gilt dann als kryptographisch sicher.
\end{bem}

\begin{bem}
	Für konventionelle Kryptoverfahren (RSA und ElGamal/DSA auf $(\FF_{p^r}^*,\cdot)$) gibt es subexponentielle Verfahren zur Lösung des DL-Problems. Dieser Vergleich schlägt sich in der Wahl der Schlüssellängen ($=$ Bitzahl der Größe des endlichen Körpers) nieder:
	Die Schlüssellänge eines elliptischen Kurven-Systems wächst etwas schneller als die dritte Wurzel der Schlüssellänge eines konventionellen Krypto-Systems mit ähnlicher kryptographischer Sicherheit: \todo{Grafik einfügen!}
\end{bem}

\begin{bem}
	\begin{itemize}
		\item Man geht davon aus, dass Kurven $E(\FF_p)$ mit $p \approx 2^{173}$. wo $\#E(\FF_p)$ einen Primteiler $\geq 2^{160}$ hat, die gleiche Sicherheit wie ein RSA-System mit $1024$ Bit bietet (für $4096$ Bit beim RSA nur etwa $313$ bei EC-System!).
		\item Durch die geringere Schlüssellänge bei Verfahren mit elliptischen Kurven kann man diese leicht auf Smart-Cards ohne Koprozessor implementieren.
		Solche Smart-Cards sind wesentlich billiger als Chip-Karten mit Koprozessor.
	\end{itemize}
\end{bem}

\begin{bem}
	Bedenken der elliptischen Kurven-Kryptographie:
	Die Nichteignung supersingulärer und anormaler Kurven kam schnell und überraschend.
	Es ist unklar, ob noch weitere ungeeignete Kurvenfamilien existieren und mit einem schnellen DL-Algorithmus angreifbar sind.
\end{bem}

%\appendix							% Anhang beginnt hier (Nummerierung mit Buchstaben)
\printindex							% Index ausgeben
\subsection*{Liste der Sätze und Definitionen}
\listtheorems{satz,falko,defn,lemma}					% listet angegebene Sätze mit Namen und Seitenzahl auf
\end{document}